\section{Materials and Methods}

\subsection*{Cow husbandry}

All procedures involving cow handling and image capturing were conducted in accordance with ethical guidelines and approved by the Virginia Tech Institutional Animal Care and Use Committee (IACUC \#22-146). The cows studied were part of the dairy herd at the Virginia Tech Dairy Complex in Blacksburg, Virginia, USA, which comprises approximately ~80\% Holstein and ~20\% Jersey cows. For the 'External' setting, the study included  100\% Holstein cows. The milking cows were housed in pens within a free-stall barn, featuring two rows of sand-bedded stalls, headlocks at the feed bunk, and two water troughs per pen. The stocking density was maintained at 100\% (i.e., one cow per stall). Heat stress was managed using automatic 48-inch diameter fans positioned over the stalls and feeding alleys. Cows were milked twice daily at 1:00 am and 12:00 pm in a double-twelve parallel milking parlor. They were fed ad libitum (with less than 5\% refusals) once daily at 8:00 am with a total mixed ration (TMR) consisting of approximately 42\% corn silage, 8\% grass hay, and 50\% concentrate on a dry matter basis. Manure from the stalls was removed at each milking session by personnel driving the cows to milking. Manure from the walking alleys within the pen was cleared two or three times daily using an automatic flushing system with recycled water. Fresh or recycled sand was added on a weekly basis.

\subsection*{Study datasets}

This study consists of two datasets: the COw LOcalization (COLO) dataset and the external dataset collected from Inner Mongolia University for model evaluation. The COLO dataset was collected using the Amazon Ring camera model Spotlight Cam Battery Pro (Ring Inc.), which offers a real-time video feed of dairy cows. Three cameras were installed in the barn: two at a height of 3.25 meters (10.66 feet) above the ground covering an area of 33.04 square meters (355.67 square feet). One camera provided a top view while the other was angled approximately 40 degrees from the horizontal to offer a side view of the cows. These are hereafter referred to as \textit{the Top-View camera} and \textit{the Side-View camera}, respectively. A third camera, termed \textit{the external camera}, was set at a lower height of 2.74 meters (9.00 feet) and covered a larger area of 77.63 square meters (835.56 square feet). Positioned 10 degrees downward from the horizontal, it captured a challenging perspective prone to occlusions among cows. Since the images collected from the third camera have a different view angle and lighting conditions, they were used as an external dataset, named \textit{ext\_VT}, for model evaluation. The second dataset, named \textit{ext\_IMU}, was collected from Inner Mongolia University, China, using surveillance cameras (Dahua cameras, M/K model) to monitor the cow daily behaviors, such as feeding, drinking, and ruminating in the barn \citep{li_new_2024}. This study sampled 250 images with 3,699 cow instances to form the \textit{ext\_IMU} dataset, that has different farm layout and camera angles from the COLO dataset. The example images from the datasets were shown in \textbf{Figure~\ref{fig:dataset}}. 

Images were captured using an Ring Application Programming Interface (API) \cite{greif_dgreifring_2024}, configured to record a ten-second video clip every 30 minutes continuously for 14 days. Since the image quality relies on the camera's internet connection, which was occasionally unstable, some images were found to be tearing or unrecognizable. Hence, the resulting dataset was manually curated for consistent quality, comprising 504 images from \textit{the Top-View camera}, 500 from \textit{the Side-View camera}, and 250 from \textit{the external camera}. These images were further categorized based on the lighting conditions: for \textit{the Top-View camera}, 296 images were captured during daylight, 118 in the evening under artificial lighting, and 90 as near-infrared images without artificial light. From \textit{the Side-View camera}, 113 images were taken in the evening, and 97 as near-infrared images. All images from \textit{the external camera} were captured during the day.

\begin{figure}[h]
    \centering
    \includegraphics[width=1\textwidth]{figure_1.jpg}
    \caption{Overview of the COLO dataset and external dataset. (a) Six instance images from the dataset with pink bounding boxes labeling the location of cows. The columns show two different view angles: top-view and side-view. The rows show three different lighting conditions: daylight, indoor, and near-infrared. (b) Example images of the external datasets: ext\_VT and ext\_IMU. (c) An visual example of the annotated image in YOLO format that is used in this study. W, H, w$_b$, and h$_b$ represent the width, height, width of the bounding box, and height of the bounding box, respectively. x and y represent the center coordinates of the bounding box.} 
    \label{fig:dataset}
\end{figure}

The image annotations were conducted using an online platform, Roboflow \cite{roboflow2023}, to define cow positions in the images. The bounding boxes were manually drawn to enclose the cow contours, providing the coordinates of the top-left corners and the width and height of the boxes. If cows were partially occluded, the invisible parts were inferred based on the adjacent visible parts. If the cow position was too far from the camera, making important body features such as the head, tail, and legs unrecognizable, the cow was excluded from the dataset. The final annotations were saved in the YOLO format \cite{ultralytics2023datasets}, where annotations were stored in a text file with one row per cow in the image, each row containing the cow's class, center coordinates, width, and height of the bounding box. The graphical representation of the annotated images is shown in \textbf{Figure~\ref{fig:dataset}c}.

\subsection*{Study 1: benchmarking model generalization across different environmental conditions}

This study aimed to examined the first hypotheses that the model performance is equally affected by the view angle and lighting conditions. To compare the performance difference between view angles and lighting conditions, we designed a Monte Carlos cross-validation strategy where models were trained on one data split (e.g., top-view angle) with sampled images, and evaluated on a fixed set of images from another split (e.g., side-view angle). There are four training configurations in this study (\textbf{Figure~\ref{fig:experiment}a}):

\begin{itemize}
    \item \textit{Baseline:} The model was trained and evaluated on the dataset characterized for all conditions, including top-view, side-view, daylight, evening, and near-infrared images. The images did not overlap between the training and test sets.
    \item \textit{Top2Side:} The model was trained on the top-view images and evaluated on the side-view images.
    \item \textit{Side2Top:} The model was trained on the side-view images and evaluated on the top-view images.
    \item \textit{Day2Night:} The model was trained on the daylight images and evaluated on the evening images, including both artificial lighting and near-infrared images.
\end{itemize}

\begin{figure}[h]
    \centering
    \includegraphics[width=1\textwidth]{figure_3.jpg}
    \caption{Experiment design for the model generalization study. (a) The experiments in the Study 1 and 2 contain four cross validation configurations to investigate the affect of different environmental conditions on model performance: Baseline, Top2Side, Side2Top, and Day2Night. (b) The study 3 compares the performance of two different initial weights for fine-tuning the model on external datasets. Noted that the size of each rectangular box is not proportional to the number of images in the dataset. It is to visually show where the images are sampled from.}
    \label{fig:experiment}
\end{figure}

To study how the training sample size affects model performance in each configuration, the testing set in the cross-validation was fixed to the same 100 images. Then, the training set size was iteratively altered from 16 to 500 images, with the sample size doubled at each step until reaching 500 images. Each training sample size was repeatedly sampled 20 times with different random seeds to ensure the robustness of the results.

To assess significant differences in mAP@0.5:0.95 performance among the four configurations, pairwise statistical tests were conducted. For each pair of configurations, data normality was first evaluated using the Shapiro-Wilk test \cite{shapiro_analysis_1965}. If both were normal, homogeneity of variances was checked with Levene's test \cite{levene_robust_1960}. Based on these tests, either a two-sample t-test (normal, equal variance) or Welch's t-test \cite{welch_generalization_1947} (normal, unequal variance) was used to compare means. If the data were not normally distributed, the Mann-Whitney U test \cite{mann_test_1947} was applied to compare the distributions. The significance level was set at p < 0.05.

\subsection*{Study 2: the correlation between model complexity and performance on the tasks of localizing cows}

To examine whether model performance improves with increasing complexity, three object detection model families were selected for this study: YOLOv11, YOLOv12, and RT-DETR. Within each family, two size variants were chosen to assess the impact of model complexity under a consistent architecture and mechanism. The selected models and their parameter counts (in parentheses) are as follows: YOLOv11n (2.6 million) and YOLOv11x (56.9 million), YOLOv12n (2.6 million) and YOLOv12x (59.1 million), and RT-DETR-l (45 million) and RT-DETR-x (86 million).

All models were evaluated using the same cross-validation configuration as in Study 1 (Figure~\ref{fig:experiment}a). For each model pair that consists of a smaller and a larger variant from the same family, performance differences were statistically tested as the same procedures as in Study 1. A significant difference between the two sizes within a family would indicate that model performance is influenced by complexity. Conversely, a lack of significant difference would suggest that comparable performance can be achieved with the smaller, more computationally efficient model.

\subsection*{Study 3: assessing the advantages of using fine-tuned model over the pre-trained model as initial model weights}

Most models are released with pre-trained weights obtained from large datasets containing millions of object instances (e.g., COCO \cite{lin2014microsoft} and ImageNet \cite{deng2009imagenet}). The pre-trained models have a general capability in recognizing common objects such as vehicles, animals, and household items. When the model is required to recognize specific objects (i.e., cows in this study), a model trained on a smaller but specific dataset is expected to have better performance. However, such advantages may not necessarily persist as the training sample size increases. Having an equally large number of samples for both the pre-trained and fine-tuned models could diminish the performance gap between the two models.

To investigate this hypothesis, we evaluated the performance of fine-tuned models on external datasets using two different initial weight configurations \textbf{(Figure~\ref{fig:experiment}b)}. To compare the impact of different initializations, we used two sets of initial weights. The first was obtained by fine-tuning a pre-trained model on the COLO dataset presented in this work (1,004 images, 9,846 cow instances), while the second came directly from a model pre-trained on the COCO dataset. With these two sets of weights, we fine-tuned models on the external datasets, \textit{ext\_VT} and \textit{ext\_IMU}, using sample sizes of 16, 128, and 200 images for fine-tuning. Evaluation was performed on a fixed set of 50 unseen images from the same external dataset source. Similar to Study 1, the training set was sampled 20 times with different random seeds to account for sampling variability. This allowed us to track the gradual performance change as more images from the new environment were incorporated. Model performance was assessed using the mAP@0.5:0.95 metric. We conducted a pairwise comparison between the two initializations for each sample size and model variant using the same statistical procedures from Study 1 and Study 2. A statistically significant difference would suggest that transferring weights from a task-specific dataset (COLO) offers a greater advantage than using general-purpose pre-trained weights from a large dataset (COCO).

\subsection*{Model training}

The model training was conducted using the Ultralytics framework \citep{ultralytics_yolo}, which provides predefined data augmentation strategies and optimized hyperparameters for the training process. Data augmentation aims to increase the diversity of the training dataset and improve the model robustness under varying conditions. For instance, horizontal flipping randomly mirrors images, helping prevent overfitting to specific orientations present in the training data. There are also other augmentation techniques that shift the color space of the images, mimicking the variations in lighting conditions. All these techniques are designed to enhance the model to better generalize to new environments that may differ from the training data. This study focused on data augmentation techniques that operate on a single image at a time, in contrast to more complex methods like CutMix \citep{yun_cutmix_2019} and MixUp \citep{zhang_mixup_2018} which combine multiple images in random, overlapping positions. This focus was maintained to minimize potential bias in evaluation introduced by the interaction of multiple images. \textbf{Figure~\ref{fig:augmentation}} shows two instance images with different applied augmentations. The complete list of data augmentation techniques used in this study is detailed in \textbf{Table~\ref{tab:augmentation}}.

\begin{table}[h]
    \caption{Data augmentation parameters of Ultralytics used in the study. The column of probability indicates the likelihood of applying the augmentation.}
    \centering
    \begin{tabular}{lp{0.5\textwidth}l}
        \toprule
        \textbf{Augmentation Parameter} & \textbf{Description} & \textbf{Probability} \\
        \midrule
        fliplr & Randomly flips the images horizontally to create a mirror image. & 0.5 \\
        \midrule
        scale & Randomly scales the images to simulate different distances between the camera and the cows & 0.5 \\
        \midrule
        translate & Randomly move the images to simulate different camera angles. & 0.1 \\
        \midrule
        hsv\_h & Randomly changes the hue of the images. & 0.015 \\
        \midrule
        hsv\_s & Randomly changes the saturation of the images. & 0.7 \\
        \midrule
        hsv\_v & Randomly changes the brightness of the images. & 0.4 \\
        \bottomrule
    \end{tabular}
    \label{tab:augmentation}
\end{table}

\begin{figure}[H]
    \centering
    \includegraphics[width=1\textwidth]{figure_2.jpg}
    \caption{Two example images with different data augmentation applied. The left image is the original image, and the middle and right images are the augmented images. The augmentations include horizontal flipping, scaling, translation, and color space changes.}
    \label{fig:augmentation}
\end{figure}

Additionally, the training process was governed by predefined hyperparameters, including the number of epochs, batch size, optimizer, and learning rate. Rather than using the number of epochs as a fixed standard, this study adopted a constant number of update steps set to 8,192 (i.e., $2^{13}$). It ensured that all experiments, regardless of dataset size, underwent the same amount of model updates. For example, if the sample size is 500, the number of epochs is 8192 / (500 / 16) = 262.4 $\approx$ 262 epochs. This approach is especially important for experiments with smaller sample sizes, where a fixed number of epochs would otherwise result in significantly fewer updates. By standardizing on update steps, smaller datasets are trained over more epochs to reach the same number of updates, ensuring a fairer comparison across experiments. Early stopping was applied with a patience of one-quarter of the total number of epochs, and training was terminated if the model’s performance, measured by mAP@0.5:0.95, did not improve within this window. The optimization process was performed using the Adam optimizer \citep{kingma_adam_2017}. A detailed list of training hyperparameters is presented in \textbf{Table~\ref{tab:hyperparameters}}.

\begin{table}[h]
    \caption{Hyperparameters for Ultralytics model training.}
    \centering
    \begin{tabular}{lp{0.5\textwidth}l}
        \toprule
        \textbf{Hyperparameter} & \textbf{Description} & \textbf{Value} \\
        \midrule
        step & Number of update steps for training & 8192 \\
        \midrule
        epochs & Number of time to iterate through the entire dataset & 8192 / ($n$ / 16)$^a$ \\
        \midrule
        batch & Number of images in each batch & 16 \\
        \midrule
        optimizer & Optimizer used for training & Adam \\
        \midrule
        lr0 & Initial learning rate & 0.01 \\
        \midrule
        lrf & Learning rate final value & 0.1 \\
        \midrule
        momentum & Momentum for the optimizer & 0.9 \\
        \bottomrule
    \end{tabular}
    \label{tab:hyperparameters}

    $^a$ The sample size $n$ is the number of images in the training dataset. The number of epochs is calculated as the total number of update steps divided the sample size divided by the batch size.

\end{table}

\subsection*{Model evaluation}

The examined YOLO models are object detection models that return positions of detected objects (i.e., cows in this study) for the evaluated images. The detections are represented by a list of bounding boxes. Regardless of specific procedures among YOLO variants for computational efficiency, such as YOLOv8, which integrates objectness scores and conditional class probabilities into a single confidence score, each detection generally consists of $4+c$ elements: the xy-coordinates, width, and height of the bounding box, and the $c$ confidence scores indicating the probability of the object belonging to each of the $c$ classes. The class with the highest confidence score is considered the predicted class of the object. To evaluate the model performance, two aspects are considered: the localization accuracy and the classification accuracy. The localization accuracy is measured by the Intersection over Union (IoU) between the predicted bounding box and the ground truth bounding box. On the other hand, the classification accuracy is measured by the precision and recall given the confidence threshold. If the confidence score of a detection is higher than the threshold, the detection is considered a positive detection. Otherwise, the detection is neglected. Mean Average Precision (mAP) is a key metric in object detection that combines both localization and classification accuracy. It is calculated by averaging the area under the precision-recall curve across all classes. This curve is generated by varying the confidence threshold from 0 to 1 for a given IoU threshold. In this study, comprehensive model performance was primarily evaluated using the averaged mAP across IoU thresholds ranging from 0.5 to 0.95, denoted as mAP@0.5:0.95. To provide a more intuitive understanding, we also report mAP at a less stringent IoU threshold of 0.5 (mAP@0.5), along with precision and recall at a confidence threshold of 0.25. The 0.25 threshold was selected as it is the default in the Ultralytics framework \citep{ultralytics_yolo}. This threshold tends to yield higher recall, which is favorable in scenarios where missing detections are more critical than false positives.

Additionally, this study examined the computational cost of different models in terms of both processing time and memory usage during training and inference. Understanding computation time is crucial for models intended for real-time applications, which typically require a minimum of 30 frames per second (FPS) to meet real-time performance standards. Memory usage is another critical factor that affects the feasibility of deploying a model on specific hardware. For instance, training a large model that demands 30 gibibytes (GiB) of video memory (VRAM) would exceed the capacity of a consumer-grade GPU like the NVIDIA 3080 Ti, which offers only 12 GiB of VRAM.

For the training phase, we measured the wall clock time per update step, which is defined as the time required to process a batch of images and update model weights, along with the peak VRAM usage based on the log files generated by the TensorBoard \citep{martin_abadi_tensorflow_2015}. During inference, FPS was measured by processing a batch of 16 images and recording the total time taken. The VRAM usage was also monitored during inference through PyTorch's build-in function \texttt{torch.cuda.max\_memory\_allocated()} \citep{paszke_pytorch_2019}. To ensure robustness, both training and inference measurements were averaged over 30 iterations. Training and inference were conducted on NVIDIA A100 80GB and NVIDIA L40S 48GB GPUs, while inference was additionally tested on NVIDIA 3080 Ti 12GB and NVIDIA 4090 24GB GPUs.

\subsection*{Software environment and hardware details}
The computing environment for this study was configured across three different platforms. As multiple models were evaluated under varying hardware and software settings, details such as software versions, GPU specifications, and other system configurations are summarized in \textbf{Table~\ref{tab:hardware}} to clearly document the differences among the computing environments.

\begin{table}[h]
    \caption{System configuration summary across three GPU setups}
    \centering
    \resizebox{\textwidth}{!}{
    \begin{tabular}{p{3cm}p{5cm}p{5cm}p{5cm}}
        \toprule
        \textbf{Category} & \textbf{System 1 (A100)} & \textbf{System 2 (L40S)} & \textbf{System 3 (3080Ti / 4090)} \\
        \midrule
        GPU & \makecell[l]{NVIDIA A100-SXM4 80GB} 
            & \makecell[l]{NVIDIA L40S 48GB} 
            & \makecell[l]{NVIDIA 4090 24GB \\ NVIDIA 3080 Ti 12GB} \\
        \midrule
            CPU / System 
            & HPE Apollo 6500  
            & Dell PowerEdge R760XA 
            & \makecell[l]{ASUS X670E + Ryzen 9 7950X} \\
        \midrule
        Operation System
            & Rocky Linux 9.5
            & Rocky Linux 9.4
            & Ubuntu 24.04.2 \\
        \midrule
        Python Version& 3.11.5 & 3.11.5 & 3.11.8 \\
        \midrule
        PyTorch Version& 2.4.0+cu121 & 2.4.0+cu121 & 2.3.1+cu121 \\
        \midrule
        NVIDIA Driver & 550.127.05 & 550.54.15 & 550.144.03 \\
        \midrule
        CUDA Version& 12.4 & 12.4 & 12.4 \\
        \midrule
        cuDNN Version& 9.0.1 & 9.0.1 & 8.9.2 \\
        \midrule
        Ultralytics Version& 8.3.89 & 8.3.89 & 8.3.96 \\
        \bottomrule
    \end{tabular}
    }
    \label{tab:hardware}
\end{table}


